Um das Transformationsverhalten der bilinearen Kovarianten unter Lorentztransformationen vollständig zu verstehen, ist es wichtig, die Rolle von \(\operatorname{det}(\Lambda)\) bei der Transformation von \(\gamma^5\) zu klären. Dies hängt mit den Paritätseigenschaften von \(\gamma^5\) zusammen.

Zunächst transformiert ein Dirac-Spinor \(\psi\) unter einer Lorentztransformation \(\Lambda\) gemäß:

\[
\psi \rightarrow S(\Lambda) \psi
\]

und der adjungierte Spinor \(\bar{\psi}\) transformiert entsprechend:

\[
\bar{\psi} \rightarrow \bar{\psi} S^{-1}(\Lambda)
\]

Betrachten wir die bilinearen Formen:

1. **Pseudoskalar \(\bar{\psi} \gamma^5 \psi\):**

Die Transformation ergibt sich aus:

\[
\bar{\psi} \gamma^5 \psi \rightarrow \bar{\psi} S^{-1}(\Lambda) \gamma^5 S(\Lambda) \psi
\]

Um die Transformationseigenschaft von \(\gamma^5\) zu verstehen, betrachten wir:

\[
S^{-1}(\Lambda) \gamma^5 S(\Lambda) = \gamma^5 \cdot \operatorname{det}(\Lambda)
\]

Dies folgt aus der Tatsache, dass \(\gamma^5\) ein Pseudoskalar ist und sich unter Paritätstransformationen mit dem Determinantenfaktor \(\operatorname{det}(\Lambda)\) ändert. Für Paritätstransformationen gilt \(\operatorname{det}(\Lambda) = -1\), da sie die Orientierung des Raumes umkehren. Daraus folgt:

\[
\bar{\psi} \gamma^5 \psi \rightarrow \operatorname{det}(\Lambda) \bar{\psi} \gamma^5 \psi
\]

Dies zeigt, dass \(\bar{\psi} \gamma^5 \psi\) als Pseudoskalar transformiert.

2. **Axialer Vektor \(\bar{\psi} \gamma^\mu \gamma^5 \psi\):**

Für diesen Ausdruck haben wir:

\[
\bar{\psi} \gamma^\mu \gamma^5 \psi \rightarrow \bar{\psi} S^{-1}(\Lambda) \gamma^\mu \gamma^5 S(\Lambda) \psi
\]

Da \(\gamma^\mu\) als Vektor transformiert, gilt:

\[
S^{-1}(\Lambda) \gamma^\mu S(\Lambda) = \Lambda_\nu^\mu \gamma^\nu
\]

Kombiniert mit der Transformationseigenschaft von \(\gamma^5\), ergibt sich:

\[
\bar{\psi} \gamma^\mu \gamma^5 \psi \rightarrow \operatorname{det}(\Lambda) \Lambda_\nu^\mu \bar{\psi} \gamma^\nu \gamma^5 \psi
\]

Dies zeigt, dass \(\bar{\psi} \gamma^\mu \gamma^5 \psi\) als axialer Vektor transformiert.

3. **Antisymmetrischer Tensor \(\bar{\psi} \sigma^{\mu \nu} \psi\):**

Für den antisymmetrischen Tensor \(\sigma^{\mu \nu} = \frac{\mathrm{i}}{2} [\gamma^\mu, \gamma^\nu]\) ergibt sich die Transformation:

\[
\bar{\psi} \sigma^{\mu \nu} \psi \rightarrow \bar{\psi} S^{-1}(\Lambda) \sigma^{\mu \nu} S(\Lambda) \psi
\]

Da \(\sigma^{\mu \nu}\) aus den \(\gamma\)-Matrizen besteht, die als Tensoren transformieren, folgt:

\[
S^{-1}(\Lambda) \sigma^{\mu \nu} S(\Lambda) = \Lambda_\alpha^\mu \Lambda_\beta^\nu \sigma^{\alpha \beta}
\]

Daraus ergibt sich:

\[
\bar{\psi} \sigma^{\mu \nu} \psi \rightarrow \Lambda_\alpha^\mu \Lambda_\beta^\nu \bar{\psi} \sigma^{\alpha \beta} \psi
\]

Dies zeigt, dass \(\bar{\psi} \sigma^{\mu \nu} \psi\) als antisymmetrischer Tensor transformiert.

%CHECK: Erklärte die Rolle von \(\operatorname{det}(\Lambda)\) bei \(\gamma^5\) Transformation, bezog sich auf Paritätseigenschaften und hergeleitet, warum \(\operatorname{det}(\Lambda) = -1\) für Paritätstransformationen gilt.